\documentclass[11pt,a4paper]{article}

\usepackage[margin=1in]{geometry}
\usepackage{amsmath,amssymb,mathtools}
\usepackage{booktabs}
\usepackage{hyperref}
\usepackage{enumitem}
\usepackage{fancyhdr}
\usepackage{titlesec}
\usepackage{xcolor}
\usepackage{setspace}
\usepackage{caption}
\usepackage{float}
\usepackage{graphicx}

\hypersetup{
  colorlinks=true,
  linkcolor=blue!60!black,
  citecolor=blue!60!black,
  urlcolor=blue!60!black
}

\titleformat{\section}{\large\bfseries}{\thesection}{1em}{}
\titleformat{\subsection}{\normalsize\bfseries}{\thesubsection}{1em}{}
\titleformat{\subsubsection}{\normalsize\itshape}{\thesubsubsection}{1em}{}

\pagestyle{fancy}
\fancyhf{}
\rhead{\small\textit{Climate Health Resilience Framework}}
\lhead{\small\textit{OZ Statistical Consulting --- World Bank}}
\cfoot{\thepage}

\setlength{\parindent}{0pt}
\setlength{\parskip}{0.8em}

\begin{document}

\begin{center}
  {\LARGE\bfseries Retrospective Climate Health Resilience Analysis}\\[0.5em]
  {\large\itshape Adapting the Burkart et al.\ (2021) Framework\\to Measure Health System Adaptation to Observed Temperature Change}\\[1.5em]
  {\normalsize Aaron Osgood-Zimmerman\\OZ Statistical Consulting\\[0.3em]
  Prepared for the World Bank\\[0.3em]
  \today}
\end{center}

\vspace{1em}
\hrule
\vspace{1.5em}

%----------------------------------------------------------------------
\section{Motivation}
%----------------------------------------------------------------------

Burkart et al.\ (2021) developed a two-part framework for estimating the global burden of disease attributable to non-optimal temperature. Their model links cause-specific mortality to daily temperature via exposure--response surfaces estimated from 64.9 million deaths across nine countries, then applies these surfaces globally to estimate attributable deaths, DALYs, and population attributable fractions (PAFs) for 204 countries from 1990--2019.

This framework is designed to answer a \textbf{forward-looking} question: \emph{How large is the temperature-attributable health burden, and how might it grow under future warming scenarios?} The existing Colombia analysis by Osorio and Pineda Lozano follows this approach, computing attributable mortality and years of potential life lost (AVPP) for 17 temperature-sensitive causes across 33 departments for 2010--2019, with projections to 2050.

We propose to \textbf{adapt this framework for a retrospective question}: \emph{How well are Colombian health systems absorbing the health burden imposed by observed temperature change?} Rather than projecting future burden, we compare:
\begin{itemize}[nosep]
  \item The burden \textbf{expected} under observed rising temperatures, assuming population vulnerability remained at a historical baseline, against
  \item The burden \textbf{actually observed}.
\end{itemize}

The gap between these two quantities is a measure of \textbf{climate health resilience}---the degree to which health systems, infrastructure, behavior, and clinical capacity have adapted to offset the additional burden imposed by warming.

%----------------------------------------------------------------------
\section{Review of the Burkart Framework}
%----------------------------------------------------------------------

\subsection{Model Structure}

The Burkart model has two stages:

\textbf{Stage 1: Exposure--response estimation.} Using individual death records from nine countries (Brazil, Chile, China, Colombia, Guatemala, Mexico, New Zealand, South Africa, USA), cause-specific mortality rates are modeled as a function of daily mean temperature and mean annual temperature zone using MR-BRT (meta-regression---Bayesian, regularized, trimmed). This produces relative risk (RR) surfaces $\text{RR}_c(T_{\text{daily}}, T_{\text{zone}})$ for each cause $c$, defined over a two-dimensional grid of daily temperatures and 23 mean annual temperature zones (6\textdegree C to 28\textdegree C).

Two distinct curve shapes emerge:
\begin{itemize}[nosep]
  \item \textbf{J-shaped} (non-external causes: IHD, stroke, COPD, LRI, CKD, diabetes, hypertensive heart disease, cardiomyopathy): mortality increases at both temperature extremes, with cold effects dominant across a wide range and heat effects concentrated at the upper tail.
  \item \textbf{Monotonically increasing} (external causes: drowning, homicide, suicide, transport injuries, disasters): mortality rises with temperature throughout.
\end{itemize}

\textbf{Stage 2: Burden attribution.} For each grid pixel, day, and cause, the model computes:
\begin{enumerate}[nosep]
  \item The \textbf{TMREL} (theoretical minimum-risk exposure level): the location- and year-specific daily temperature associated with the lowest combined mortality across all included causes.
  \item The \textbf{PAF}: the fraction of deaths attributable to temperature above (heat) or below (cold) the TMREL.
  \item The \textbf{attributable burden}: total deaths $\times$ PAF, aggregated to the desired spatial and temporal resolution.
\end{enumerate}

\subsection{Current Colombia Implementation}

In the existing pipeline (Script~11, \texttt{11\_carga\_atribuible.R}), the attributable burden for department $d$, year $t$, cause $c$ is:
\begin{equation}
  B_{\text{obs}}(t,d,c) \;=\; M_{\text{obs}}(t,d,c) \;\times\; \text{PAF}\bigl(T_{\text{obs}}(t,d)\bigr) \;\times\; \text{SEV}\bigl(T_{\text{obs}}(t,d)\bigr)
  \label{eq:current}
\end{equation}
where:
\begin{itemize}[nosep]
  \item $M_{\text{obs}}(t,d,c)$ = observed deaths from cause $c$ in department $d$, year $t$ (aggregated from daily, individual-level DANE vital registration records)
  \item $\text{PAF}(\cdot)$ = population-attributable fraction computed from the Burkart RR curves relative to the TMREL, weighted across pixels by population share
  \item $\text{SEV}(\cdot)$ = summary exposure value, a risk-weighted prevalence of temperature exposure ranging from 0 (no risk exposure) to 1 (maximum possible risk)
\end{itemize}

Key features of the current implementation:
\begin{itemize}[nosep]
  \item The \textbf{RR curves are fixed}---estimated from pooled multi-country data by Burkart et al.\ and do not vary over time.
  \item The \textbf{TMREL is averaged} across 2010--2019, treated as constant over the study period.
  \item What varies annually: \textbf{observed temperatures} (affecting PAF and SEV) and \textbf{observed mortality counts}.
  \item PAFs are computed at the pixel-day level, then aggregated to department-day by population weighting.
  \item Heat and cold PAFs are computed separately (temperatures above vs.\ below the TMREL).
\end{itemize}

%----------------------------------------------------------------------
\section{Proposed Retrospective Resilience Framework}
%----------------------------------------------------------------------

\subsection{Core Idea}

The attributable burden in any year is the product of three time-varying components:
\begin{equation}
  B(t) \;=\; \underbrace{M(t)}_{\text{vulnerability}} \;\times\; \underbrace{\text{PAF}\bigl(T(t)\bigr) \;\times\; \text{SEV}\bigl(T(t)\bigr)}_{\text{exposure}}
  \label{eq:decomp}
\end{equation}

Changes in $B(t)$ over time arise from changes in \emph{any} of these components. To measure health system resilience, we need to disentangle \textbf{how much of the burden change is driven by shifting temperatures} (exposure) versus \textbf{how much is offset by reduced vulnerability} (adaptation).

\subsection{Baseline Period}

Define a baseline period $t_0$ (e.g., 2010--2012). From this period, compute and freeze:
\begin{itemize}[nosep]
  \item $m_0(d,c,a,s)$: baseline cause-specific mortality \emph{rates} (deaths per 100,000) by department $d$, cause $c$, age group $a$, and sex $s$.
  \item $\text{PAF}_0(d,c)$: PAFs computed from the baseline-period daily temperature distributions.
  \item $\text{SEV}_0(d,c)$: SEVs computed from the baseline-period temperature distributions.
  \item $T_0(d)$: the daily temperature distribution during the baseline period.
\end{itemize}

\subsection{Counterfactual Scenarios}

We construct four burden estimates for each year $t$, department $d$, and cause $c$:

\medskip

\begin{table}[H]
\centering
\small
\caption{Counterfactual burden scenarios}
\label{tab:scenarios}
\begin{tabular}{@{}llll@{}}
\toprule
\textbf{Scenario} & \textbf{Mortality} & \textbf{Temperature} & \textbf{Interpretation} \\
\midrule
$B_{\text{obs}}(t)$ & Observed $M(t)$ & Observed $T(t)$ & What actually happened \\[0.3em]
$B_{\text{noWarm}}(t)$ & Observed $M(t)$ & Baseline $T_0$ & What if temperatures hadn't changed? \\[0.3em]
$B_{\text{noAdapt}}(t)$ & Baseline rates $\times$ $\text{Pop}(t)$ & Observed $T(t)$ & What if vulnerability hadn't changed? \\[0.3em]
$B_{\text{base}}(t)$ & Baseline rates $\times$ $\text{Pop}(t)$ & Baseline $T_0$ & Pure demographic growth effect \\
\bottomrule
\end{tabular}
\end{table}

\medskip

Formally:
\begin{align}
  B_{\text{obs}}(t,d,c) &= M_{\text{obs}}(t,d,c) \times \text{PAF}\bigl(T_{\text{obs}}(t,d)\bigr) \times \text{SEV}\bigl(T_{\text{obs}}(t,d)\bigr) \label{eq:obs}\\[0.5em]
  B_{\text{noWarm}}(t,d,c) &= M_{\text{obs}}(t,d,c) \times \text{PAF}_0(d) \times \text{SEV}_0(d) \label{eq:nowarm}\\[0.5em]
  B_{\text{noAdapt}}(t,d,c) &= m_0(d,c) \times \text{Pop}(t,d) \times \text{PAF}\bigl(T_{\text{obs}}(t,d)\bigr) \times \text{SEV}\bigl(T_{\text{obs}}(t,d)\bigr) \label{eq:noadapt}\\[0.5em]
  B_{\text{base}}(t,d,c) &= m_0(d,c) \times \text{Pop}(t,d) \times \text{PAF}_0(d) \times \text{SEV}_0(d) \label{eq:base}
\end{align}

where $m_0(d,c)$ denotes baseline-period mortality \emph{rates} and $\text{Pop}(t,d)$ is the actual population in year $t$.

\subsection{Derived Quantities}

\subsubsection{Excess Burden from Warming}

The additional burden attributable to observed temperature change, holding vulnerability at its actual (improving) level:
\begin{equation}
  \Delta B_{\text{warming}}(t) = B_{\text{obs}}(t) - B_{\text{noWarm}}(t)
  \label{eq:excess}
\end{equation}

If temperatures have shifted toward more days above the TMREL, this quantity will be positive, indicating that warming is imposing additional health burden.

\subsubsection{Burden Absorbed by Adaptation}

The burden that \emph{would have occurred} under current temperatures if vulnerability had remained at baseline, minus what actually occurred:
\begin{equation}
  \Delta B_{\text{adapted}}(t) = B_{\text{noAdapt}}(t) - B_{\text{obs}}(t)
  \label{eq:adapted}
\end{equation}

When positive, this quantity reflects the burden successfully absorbed by improvements in health systems, clinical capacity, infrastructure, and behavior.

\subsubsection{Climate Health Resilience Index}

We define a normalized resilience index:
\begin{equation}
  \boxed{\;R(t,d,c) = 1 - \frac{B_{\text{obs}}(t,d,c)}{B_{\text{noAdapt}}(t,d,c)}\;}
  \label{eq:resilience}
\end{equation}

Interpretation:
\begin{itemize}[nosep]
  \item $R = 0$: No adaptation. Actual burden matches what baseline vulnerability would predict under current temperatures.
  \item $R > 0$: Positive adaptation. The health system is absorbing some of the temperature-attributable burden. $R = 1$ would imply complete absorption (zero observed attributable burden).
  \item $R < 0$: Maladaptation. Actual burden \emph{exceeds} what baseline vulnerability would predict, suggesting increased fragility.
\end{itemize}

This index can be computed at any aggregation level: nationally, by department, by cause, by age group, or by heat vs.\ cold.

\subsubsection{Full Decomposition}

The total change in attributable burden from the baseline to year $t$ can be decomposed additively using a Das Gupta-style approach:
\begin{equation}
  \Delta B(t) = B_{\text{obs}}(t) - B_{\text{base}}(t_0) = \Delta B_{\text{temperature}} + \Delta B_{\text{vulnerability}} + \Delta B_{\text{population}} + \Delta B_{\text{interaction}}
  \label{eq:dasgupta}
\end{equation}

where:
\begin{align*}
  \Delta B_{\text{temperature}} &= \text{change in burden due to shifting temperature distributions} \\
  \Delta B_{\text{vulnerability}} &= \text{change due to shifting cause-specific mortality rates} \\
  \Delta B_{\text{population}} &= \text{change due to population growth and aging} \\
  \Delta B_{\text{interaction}} &= \text{residual interaction among the above}
\end{align*}

This decomposition follows standard demographic methods for decomposing rates that are products of multiple factors. It reveals whether burden is rising primarily because of warming, or whether health improvements are large enough to offset the temperature signal.

%----------------------------------------------------------------------
\section{Interpreting ``Adaptation''}
%----------------------------------------------------------------------

The gap between $B_{\text{noAdapt}}$ and $B_{\text{obs}}$ captures \textbf{everything that changed about population vulnerability}, including:

\begin{enumerate}[nosep]
  \item \textbf{Temperature-specific adaptation}: early warning systems, cooling centers, air conditioning penetration, urban greening, behavioral changes during heat events.
  \item \textbf{General health system improvement}: better cardiovascular care, stroke treatment units, improved diabetes management---interventions that reduce underlying mortality regardless of temperature.
  \item \textbf{Epidemiological transition}: shifts in cause-of-death composition independent of climate.
\end{enumerate}

With 10 years of Colombian data, cleanly separating these mechanisms is difficult. However, two strategies provide partial identification:

\textbf{Cross-cause comparison.} Compute $R(t,d,c)$ separately for each cause. If the resilience index is high for IHD but near zero for drowning, cardiovascular care improvements---not temperature-specific interventions---are likely driving the adaptation signal. Causes where $R$ remains low despite strong temperature trends are candidates for targeted climate adaptation investment.

\textbf{Cross-department comparison.} Compare departments with similar temperature trends but different health system capacity (measured by, e.g., hospital bed density, health expenditure per capita, or urbanization). Departments with stronger health systems should show higher $R$ values, isolating the health-system-specific contribution.

%----------------------------------------------------------------------
\section{Implementation Plan}
%----------------------------------------------------------------------

The proposed analysis requires modifications primarily to \textbf{Script~11} (\texttt{11\_carga\_atribuible.R}). All upstream scripts (mortality cleaning, temperature processing, exposure--response curve preparation) remain unchanged.

\subsection{Step 1: Compute Baseline Parameters}

Using data from 2010--2012:
\begin{itemize}[nosep]
  \item Compute cause-specific mortality \emph{rates} by department, age group, and sex: $m_0(d,c,a,s)$.
  \item Compute $\text{PAF}_0(d,c)$ and $\text{SEV}_0(d,c)$ from the 2010--2012 daily temperature distributions using the existing Burkart RR curves and TMREL values.
  \item Store the baseline daily temperature distribution $T_0(d)$ for the ``no warming'' counterfactual.
\end{itemize}

\subsection{Step 2: Compute All Four Burden Scenarios}

For each year $t \in \{2010, \ldots, 2019\}$, compute:
\begin{enumerate}[nosep]
  \item $B_{\text{obs}}(t)$ --- the existing calculation (already implemented).
  \item $B_{\text{noWarm}}(t)$ --- substitute $\text{PAF}_0$ and $\text{SEV}_0$ in place of year-$t$ values; use actual mortality.
  \item $B_{\text{noAdapt}}(t)$ --- substitute $m_0 \times \text{Pop}(t)$ in place of actual mortality; use actual temperature-derived PAFs and SEVs.
  \item $B_{\text{base}}(t)$ --- substitute both baseline mortality rates and baseline PAFs/SEVs.
\end{enumerate}

\subsection{Step 3: Compute Derived Quantities}

\begin{itemize}[nosep]
  \item Excess burden from warming: $\Delta B_{\text{warming}}(t)$ per Equation~\eqref{eq:excess}.
  \item Burden absorbed by adaptation: $\Delta B_{\text{adapted}}(t)$ per Equation~\eqref{eq:adapted}.
  \item Resilience index: $R(t,d,c)$ per Equation~\eqref{eq:resilience}.
  \item Full Das Gupta decomposition per Equation~\eqref{eq:dasgupta}.
\end{itemize}

\subsection{Step 4: Visualization and Output}

New outputs for a modified Script~12:
\begin{itemize}[nosep]
  \item \textbf{Time series panels}: $B_{\text{obs}}$, $B_{\text{noAdapt}}$, and $B_{\text{noWarm}}$ by department, with the adaptation gap shaded.
  \item \textbf{Choropleth maps}: resilience index $R$ by department, nationally and by cause group.
  \item \textbf{Decomposition bar charts}: stacked contributions of temperature, vulnerability, population, and interaction to total burden change, by department or cause.
  \item \textbf{Cause-level heatmaps}: $R(t,d,c)$ for all 17 causes $\times$ 33 departments, identifying where adaptation is weakest.
\end{itemize}

%----------------------------------------------------------------------
\section{Contribution and Policy Relevance}
%----------------------------------------------------------------------

This reframing shifts the analytical question from \emph{``How bad will climate change be?''} to \textbf{\emph{``How well are health systems coping with climate change that has already occurred?''}} This is directly actionable for the World Bank's health system strengthening agenda:

\begin{itemize}[nosep]
  \item \textbf{Identifies adaptation gaps}: departments and causes where the resilience index is lowest are priority targets for climate-health investment.
  \item \textbf{Quantifies adaptation value}: the total $\Delta B_{\text{adapted}}$ across Colombia provides a monetary and DALY-denominated estimate of what health system improvements have already ``bought'' in terms of climate resilience.
  \item \textbf{Distinguishes temperature-specific vs.\ general health gains}: cross-cause analysis reveals whether resilience is driven by broad health improvements or targeted climate adaptation, informing whether future investment should target climate-specific interventions (early warning systems, cooling infrastructure) or general health system capacity.
  \item \textbf{Uses only observed data}: unlike projection-based analyses, this approach requires no climate model outputs, emissions scenarios, or assumptions about future adaptation---only vital registration, temperature records, and population data that Colombia already collects.
\end{itemize}

%----------------------------------------------------------------------
\section*{References}
%----------------------------------------------------------------------

\begin{enumerate}[nosep, label={[\arabic*]}]
  \item Burkart KG, Brauer M, Aravkin AY, et al.\ Estimating the cause-specific relative risks of non-optimal temperature on daily mortality: a two-part modelling approach applied to the Global Burden of Disease Study.\ \textit{Lancet}. 2021;398:685--697.
  \item Murray CJL, Lopez AD.\ On the comparable quantification of health risks: lessons from the Global Burden of Disease Study.\ \textit{Epidemiology}. 1999;10:594--605.
  \item Das Gupta P.\ Standardization and decomposition of rates: a user's manual.\ \textit{U.S.\ Bureau of the Census, Current Population Reports}. 1993;P23--186.
\end{enumerate}

\end{document}
